% Capítulo 4: Marco Metodológico

\chapter{Marco Metodológico}

\label{Chapter4}

\lhead{Capítulo 4. \emph{Marco Metodológico}}

\textit{[En un párrafo explica brevemente los contenidos de este capítulo antes de la distribución en distintos apartados. Describe el método que has seguido para desarrollar el proyecto. Este apartado debe contener las principales características de la investigación desarrollada. Se debe describir el proceso que implica la detección de la necesidad o interés del trabajo, la búsqueda de la información y las fuentes documentales empleadas, el tipo de investigación (cualitativa, cuantitativa o sociocrítica; de tipo descriptiva, observación, analítica, estudio de casos, intervención, etnográfica, comparada, proyecto, investigación-acción, buena práctica, etc.). Los apartados a considerar ---aunque puedes variarlo para adecuarlo a tu proyecto concreto, ampliando o reduciendo apartados--- son los siguientes.]}

\textit{[Se recomienda un tamaño de 5 a 15 páginas para este capítulo.]}

%----------------------------------------------------------------------------------------

\section{Población y muestra}

\textit{[Describe la población objeto de estudio y el muestreo realizado para obtener la muestra que representa dicha población.]}

%----------------------------------------------------------------------------------------

\section{Objetivos de investigación}

\textit{[¿Qué pretendes evaluar? Si los objetivos son los mismos que los del Capítulo 2 no hace falta que incluyas este apartado.]}

%----------------------------------------------------------------------------------------

\section{Variables intervinientes e instrumentos de evaluación}

\textit{[Describe qué vas a medir y a través de qué instrumentos.]}

%----------------------------------------------------------------------------------------

\section{Diseño experimental}

\textit{[Describe el método de investigación elegido y su procedimiento.]}

%----------------------------------------------------------------------------------------

\section{Modelo de análisis de datos}

\textit{[Describe cómo vas a analizar los datos que obtengas en tu investigación: si los vas a analizar cuantitativa o cualitativamente, si vas a realizar análisis estadísticos y mediante qué programas informáticos, etc.]}
